% Generated draft: requires mathematical and source review.
\documentclass[11pt]{article}
\usepackage[margin=25mm]{geometry}
\usepackage{fontspec}
\setmainfont{Latin Modern Roman}
\usepackage{amsmath,amssymb,mathtools,unicode-math}
\setmathfont{Latin Modern Math}
\usepackage{enumitem,needspace,xurl,hyperref}
\setlength{\parindent}{0pt}
\setlength{\parskip}{5pt}
\setlength{\emergencystretch}{3em}
\newcommand{\N}{\mathbb{N}}
\newcommand{\Z}{\mathbb{Z}}
\newcommand{\Q}{\mathbb{Q}}
\newcommand{\R}{\mathbb{R}}
\newcommand{\C}{\mathbb{C}}
\newcommand{\F}{\mathbb{F}}
\newcommand{\A}{\mathbb{A}}
\newcommand{\PP}{\mathbb{P}}
\newcommand{\E}{\mathbb{E}}
\newcommand{\eps}{\varepsilon}
\newcommand{\dd}{\,\mathrm{d}}
\begin{document}
\section*{UnsolvedMath ID: 30006335}
\textbf{Source problem code:} catalogue entry
\subsection{Definitions and mathematical statement}
Write \(M_i:=H^i(\widehat U)\), retaining the notation of the supplied record. The extracted question is whether, for the data and cohomological degrees intended in the source,
\[
\forall\alpha\in M_i,\qquad
\alpha\text{ is strongly }p\text{-divisible}
\ \Longrightarrow\ 
\alpha\text{ is torsion}.
\]
The supplied background describes \(M_i\) as completed de Rham cohomology. However, the record does not specify the geometric object \(U\), its base, the meaning of \(\widehat U\), the completion used in cohomology, the coefficient ring, the allowed values of \(p\) and \(i\), or the definitions of strong \(p\)-divisibility and torsion. Consequently, the displayed implication preserves the extracted question but is not a self-contained mathematical formulation. In particular, strong \(p\)-divisibility must not silently be replaced by ordinary divisibility by every power of \(p\), and the intended meaning of torsion must be recovered from the source.

The supplied assessment also reports non-torsion strongly \(p\)-divisible classes in a related uncompleted cohomology group. That group and its relationship to \(H^i(\widehat U)\) are not specified here; this reported phenomenon is background, not a counterexample to the displayed question.
\subsection{Short English statement}
The question asks whether every strongly \(p\)-divisible class in the intended completed de Rham cohomology is torsion. Essential definitions and hypotheses are missing from the supplied extraction, so the statement requires source consultation and human mathematical review.
\subsection{Sources}
\begin{itemize}
\item Local registry: \texttt{data/unsolvedmath/problems.json}, supplied record with permanent UnsolvedMath ID \(30006335\), source problem code \texttt{catalogue entry}, titled ``Torsion of Strongly p-Divisible Cohomology Classes.''
\item Supplied report citation: \emph{Non-Archimedean Geometry and Applications} (2025), Oberwolfach Reports, DOI 10.4171/owr/2025/29, pp.\ 1517--1578. Supplied source URL and literature reference labeled ``OWR report'': \url{https://doi.org/10.4171/owr/2025/29}. No individual proposer or precise problem-page locator is supplied.
\item Literature assessment dated exactly \texttt{2026-08-22}: the supplied curation labels the problem open, calls it the \(p\)-GGHS problem, and reports non-torsion strongly \(p\)-divisible classes in a related uncompleted cohomology group. This is supplied dated curation, not an independently verified assessment; its time-sensitive status and cited claims require independent verification.
\item The registry labels the statement \texttt{corrected\_verified}, with the qualification that only formula delimiters, notation, grammar, and typography were repaired and that mathematical content and scope were preserved. This editorial label does not establish mathematical correctness or supply the missing definitions.
\end{itemize}
\end{document}
